% !PRINTED 272
\weberpage{272}
We assume that we can in some way construct a series of $m + 1$ continuous functions
\weq{1}{f(x), f_1(x), f_2(x) \ldots f_m(x), \quad (\mathfrak{R})}
possessing the following properties:

1. Of the functions $f_\nu$, no two consecutive ones shall vanish simultaneously in the interval $\Delta$.

2. The last of them, $f_m(x)$, shall not vanish at all in the interval $\Delta$, and hence shall keep an unchanging sign.

3. If a middle term, say $f_\nu(x)$, vanishes for some $x$ in the interval $\Delta$, then for this $x$ the two adjacent functions $f_{\nu-1}(x)$ and $f_{\nu+1}(x)$ shall have opposite signs.

4. If $f(x)$ vanishes in the interval $\Delta$, then $f_1(x)$ shall have the same sign for this value of $x$ as $f'(x)$.

Such a series of functions ($\mathfrak{R}$) we shall call a \gesperrt{Sturm chain}. How one can form it we shall see later; first, consequences are to be drawn from the definition.

For every value of $x$ for which none of the functions of ($\mathfrak{R}$) vanishes, each of these functions has a definite sign. We count a sign change whenever, in running through the chain from left to right, a positive term is followed by a negative one, or a negative term by a positive one. The number of sign changes (variations) counted in this way for a given value of $x$ we shall denote by $V(x)$. If a middle term $f_\nu(x)$ vanishes, then by 3. we are to count one sign change in passing from $f_{\nu-1}(x)$ to $f_{\nu+1}(x)$.

If $\alpha$, $\beta$ are two values for which none of the functions $f_\nu$ vanishes, and if we now let $x$ increase continuously from $\alpha$ to $\beta$, then a change in the number of sign changes can occur only when one of the functions of ($\mathfrak{R}$) changes its sign, and hence passes through the value zero (because of the assumed continuity).

But if this is a middle function, the number of sign changes does not alter (by 3). For in
\[
f_{\nu-1}, \; f_\nu, \; f_{\nu+1}
\]
%CONT>>
